\documentclass{article}
\usepackage{algorithm, algorithmic, amsfonts, amsmath}

\begin{document}
\begin{algorithm}[!h]
\renewcommand{\algorithmicrequire}{\textbf{Input:}}
\renewcommand{\algorithmicensure}{\textbf{Output:}}
	\caption{PGAN: learning a binary classifier from a PU dataset with a GAN}
	\begin{algorithmic}[1]
	    \REQUIRE
	    Positive set $P$; unlabeled set $U$ (mixture of positives and negatives);
	    noise prior $p_{z}(z)$; batch size $m$
	    \ENSURE
	    Binary classifier $D_{B}$
	    \STATE \textbf{Stage 1: adversarial training on the unlabeled set}
	    \FOR{each GAN epoch, each batch}
	    \STATE Sample $\{x_{U}^{(i)}\}_{i=1}^{m}\sim p(x_{U})$ and $\{z^{(i)}\}_{i=1}^{m}\sim p_{z}$;
	    set $x_{F}^{(i)}=G(z^{(i)})$
	    \STATE Update $D_{U}$ by minimising
	    $-\frac{1}{m}\sum_{i}\left[\log D_{U}(x_{U}^{(i)})+\log\!\left(1-D_{U}(x_{F}^{(i)})\right)\right]$
	    \STATE Update $G$ by minimising
	    $-\frac{1}{m}\sum_{i}\log D_{U}(G(z^{(i)}))$ % the generator competes with the unlabeled set only and never sees P
	    \ENDFOR
	    \STATE Freeze $G$
	    \STATE \textbf{Stage 2: train the classifier with generated counter-examples}
	    \FOR{each classifier epoch, each batch}
	    \STATE Sample $\{x_{P}^{(i)}\}\sim p(x_{P})$ labelled $1$, and
	    $x_{F}^{(i)}=G(z^{(i)})$ labelled $0$ % generated samples stand in as negatives, which sidesteps selecting negatives from U
	    \STATE Update $D_{B}$ by minimising
	    $-\frac{1}{m}\sum_{i}\left[\log D_{B}(x_{P}^{(i)})[1]+\log D_{B}(x_{F}^{(i)})[0]\right]$
	    \ENDFOR
		\RETURN $D_{B}$;\quad predict by $\arg\max D_{B}(x)$
	\end{algorithmic}
\end{algorithm}
% Chiaroni, Rahal, Hueber, Dufaux, ICIP 2018, "Learning with a Generative Adversarial
% Network from a Positive Unlabeled Dataset for Image Classification".
% Position: rather than selecting reliable negatives from U, the method generates negative substitutes with a GAN,
% which is why the survey files it under branch G, generative and distribution learning, rather than branch A.
%
% Point needing verification: this pseudocode was written from secondary material, in which lines 3 to 5 describe G as
% competing against the whole unlabeled set. But U is a mixture of both classes, so if G converges to that mixture the
% generated samples must contain positives, and using them as negatives on line 10 introduces a systematic bias.
% The group's later work, titled "Generating Relevant Counter-Examples from a Positive Unlabeled Dataset", indicates
% that the generation target should be the negative component of U rather than the whole mixture.
% How the original achieves this, whether by early stopping, by subtracting distributions using P, or through another
% loss, is not stated in the material at hand and must be checked against the original. It decides whether the method holds.
\end{document}
